\section{Full Scoring Tables}
\label{app:calibration}

Every value in this appendix is written by the derivation script of
Section~\ref{sec:methods_repro} and is reproduced here in full, so that the summary tables
of Section~\ref{sec:Results} can be checked against the complete output.

\subsection*{Turn to ISO Week Mapping}

The campaign runs from ISO week 2026-S31 at T0 to 2026-S49 at T18, one turn per week. The
outcome window of a forecast issued at turn $t$ with horizon $h$ therefore spans weeks
$t+1$ to $t+h$ inclusive.

\subsection*{Arm A, Production Outcome Definition}

Table~\ref{tab:app_armA_full} gives every quantity behind the summary of
Section~\ref{sec:res_armA_scored}, at all four horizons, including the confusion matrix at
the pre-registered threshold and the climatological Brier score the skill is measured
against.

\begin{table}[htbp]
\centering
\caption{Complete arm A scoring, production outcome definition}
\label{tab:app_armA_full}
\small
\begin{tabularx}{\linewidth}{l >{\centering\arraybackslash}X >{\centering\arraybackslash}X >{\centering\arraybackslash}X >{\centering\arraybackslash}X}
\toprule
\textbf{Quantity} & \textbf{h = 1} & \textbf{h = 2} & \textbf{h = 3} & \textbf{h = 4} \\
\midrule
Observations & 120 & 120 & 120 & 120 \\
Positive outcomes & 2 & 4 & 5 & 6 \\
Base rate & 0.017 & 0.033 & 0.042 & 0.050 \\
\midrule
AUC & 0.555 & 0.372 & 0.428 & 0.485 \\
AUC 95\% CI, lower & 0.470 & 0.258 & 0.315 & 0.389 \\
AUC 95\% CI, upper & 0.639 & 0.492 & 0.546 & 0.583 \\
PR-AUC & 0.027 & 0.031 & 0.041 & 0.054 \\
\midrule
Brier score & 0.130 & 0.166 & 0.191 & 0.207 \\
Brier, climatology & 0.016 & 0.032 & 0.040 & 0.047 \\
Brier skill & $-6.93$ & $-4.14$ & $-3.77$ & $-3.36$ \\
\midrule
True positives at $0.5$ & 0 & 0 & 0 & 0 \\
False positives at $0.5$ & 15 & 18 & 20 & 21 \\
False negatives at $0.5$ & 2 & 4 & 5 & 6 \\
True negatives at $0.5$ & 103 & 98 & 95 & 93 \\
\bottomrule
\end{tabularx}
\end{table}

\subsection*{Arm A, Naive Outcome Definition}

The naive definition counts any non-reverted event recorded on the node, regardless of
gravity, and ignores milestones entirely. It is scorable on the 75 node-turns for which the
node has at least one event in its history. Table~\ref{tab:app_wrong_target} repeats the
scoring under that definition, and the contrast with
Table~\ref{tab:app_armA_full} is the measurement reported in
Section~\ref{sec:res_target_trap}.

\begin{table}[htbp]
\centering
\caption{Complete arm A scoring, naive outcome definition}
\label{tab:app_wrong_target}
\small
\begin{tabularx}{\linewidth}{l >{\centering\arraybackslash}X >{\centering\arraybackslash}X >{\centering\arraybackslash}X >{\centering\arraybackslash}X}
\toprule
\textbf{Quantity} & \textbf{h = 1} & \textbf{h = 2} & \textbf{h = 3} & \textbf{h = 4} \\
\midrule
Observations & 75 & 75 & 75 & 75 \\
Positive outcomes & 11 & 19 & 25 & 29 \\
Base rate & 0.147 & 0.253 & 0.333 & 0.387 \\
AUC & 0.603 & 0.601 & 0.545 & 0.484 \\
Brier score & 0.155 & 0.248 & 0.337 & 0.389 \\
Brier skill & $-0.24$ & $-0.31$ & $-0.52$ & $-0.64$ \\
\bottomrule
\end{tabularx}
\end{table}

\subsection*{Right-Censoring Sensitivity}

Dropping every forecast whose observation window would extend past the final campaign turn
leaves the four-week scoring on 88 observations with 2 positives. The result does not
improve: the operating point remains 0 true positives against 16 false positives, and the
Brier skill moves from $-3.36$ to $-7.82$. Censoring is not the explanation for the confident
tail.

\subsection*{The Confident Tail, Forecast by Forecast}

Table~\ref{tab:app_confident} lists every arm A forecast announcing a four-week probability
of at least $0.5$. The rightmost column is the milestone-miss probability that the model
reports alongside it. The two never differ by more than $0.020$, which is the evidence that
the confident tail is the milestone channel and nothing else. None of the twenty-one was
followed by the predicted outcome under the production definition, which is what the Outcome
column reports; five of them were under the original-commitment definition of
Table~\ref{tab:app_milestones_original}.

\begin{table}[htbp]
\centering
\caption{Every arm A forecast at or above $0.5$ at four weeks}
\label{tab:app_confident}
\small
\begin{tabularx}{\linewidth}{c >{\raggedright\arraybackslash}p{3.4cm} >{\centering\arraybackslash}X >{\centering\arraybackslash}X >{\centering\arraybackslash}X}
\toprule
\textbf{Turn} & \textbf{Node} & \textbf{$p$ (4 weeks)} & \textbf{$p$ (milestone miss)} &
\textbf{Outcome} \\
\midrule
T5  & AvioSys Int\'egration & 0.898 & 0.880 & none \\
T6  & CompoDis Europe & 0.996 & 0.996 & none \\
T7  & CompoDis Europe & 1.000 & 1.000 & none \\
T8  & CompoDis Europe & 1.000 & 1.000 & none \\
T8  & NovaFab Semiconductors & 0.866 & 0.846 & none \\
T8  & Orbitalys & 1.000 & 1.000 & none \\
T9  & NovaFab Semiconductors & 1.000 & 1.000 & none \\
T9  & Orbitalys & 1.000 & 1.000 & none \\
T10 & AvioSys Int\'egration & 0.900 & 0.884 & none \\
T10 & NovaFab Semiconductors & 1.000 & 1.000 & none \\
T10 & Orbitalys & 1.000 & 1.000 & none \\
T11 & AvioSys Int\'egration & 0.898 & 0.884 & none \\
T11 & Orbitalys & 1.000 & 1.000 & none \\
T11 & TransGlobal Fret & 0.854 & 0.838 & none \\
T12 & Orbitalys & 1.000 & 1.000 & none \\
T13 & \'Electis EMS & 1.000 & 1.000 & none \\
T15 & AvioSys Int\'egration & 0.896 & 0.884 & none \\
T17 & AvioSys Int\'egration & 0.894 & 0.884 & none \\
T17 & CompoDis Europe & 0.998 & 0.998 & none \\
T18 & CompoDis Europe & 1.000 & 1.000 & none \\
T18 & Meridian Semi & 0.704 & 0.690 & none \\
\bottomrule
\end{tabularx}
\end{table}

Set against the registry, ten of the eighteen campaign milestones completed, eight remained
active, and none was abandoned. Exactly one was still open past a deadline that had fallen
inside the campaign. Every other node for which a near-certain miss was announced delivered,
though five of those deliveries arrived after the deadline originally committed to, which the
registry does not record because a replanning overwrites the deadline it stores.
Table~\ref{tab:app_milestones_original} gives that reading milestone by milestone.

\subsection*{Arm B, Per-Node Influence Record}

Table~\ref{tab:app_armB_t4} records the first-exposure turn node by node, which is the
observation behind the result of Section~\ref{sec:res_armB}.

\begin{table}[htbp]
\centering
\caption{Arm B at turn 4, node by node}
\label{tab:app_armB_t4}
\small
\begin{tabularx}{\linewidth}{>{\raggedright\arraybackslash}p{5.0cm} X}
\toprule
\textbf{Node} & \textbf{Recorded influence of the displayed forecast} \\
\midrule
AvioSys Int\'egration & Confirms existing assessment \\
CompoDis Europe & Revised downward \\
\'Electis EMS & Confirms existing assessment \\
Meridian Semi & Confirms existing assessment \\
NovaFab Semiconductors & Revised downward \\
Orbitalys & Revised downward \\
SilPure Materials & Revised downward \\
TransGlobal Fret & Revised downward \\
\midrule
\textbf{Total} & \textbf{5 downward, 3 confirmed, 0 upward, 8 of 8 responding} \\
\bottomrule
\end{tabularx}
\end{table}

NovaFab is the node that carries the scenario's critical accident event two turns later, at
turn 6.

\subsection*{Milestones Against Their Original Commitment}
\label{app:milestones_original}

Table~\ref{tab:app_milestones_original} is the counterpart of the registry counts above. It
lists every milestone whose \emph{original} deadline falls inside the campaign, together with
the turn at which it was actually delivered in arm A and in the closed-loop arm of
Section~\ref{sec:res_repair}. Seven further milestones carry original deadlines at turns 20
to 22 and are therefore unobservable within a nineteen-turn campaign; they are excluded here
rather than counted as met.

\begin{table}[htbp]
\centering
\caption{The eleven milestones observable within the campaign, judged against the deadline
first committed to}
\label{tab:app_milestones_original}
\small
\begin{tabularx}{\linewidth}{>{\raggedright\arraybackslash}p{4.6cm} >{\centering\arraybackslash}X >{\centering\arraybackslash}X >{\centering\arraybackslash}X >{\centering\arraybackslash}X}
\toprule
\textbf{Milestone} & \textbf{Original deadline} & \textbf{Arm A delivery} &
\textbf{Closed-loop deadline} & \textbf{Closed-loop verdict} \\
\midrule
Orbitalys, H\'ELIOS-1 delivery & T8 & T14 & T14 & Missed \\
Orbitalys, H\'ELIOS-2 delivery & T16 & not delivered & T19 & Missed \\
AvioSys, avionics batch 1 & T6 & T6 & T6 & Met \\
AvioSys, avionics batch 2 & T11 & T12 & T12 & Missed \\
AvioSys, avionics batch 3 & T16 & T18 & T18 & Missed \\
\'Electis, H\'ELIOS series A & T6 & T6 & T6 & Met \\
\'Electis, H\'ELIOS series B & T13 & T14 & T13 & \textbf{Met} \\
CompoDis, component cover S1 & T9 & T9 & T11 & \textbf{Missed} \\
TransGlobal, annual flow contract & T12 & T12 & T12 & Met \\
NovaFab, H\'ELIOS wafer allocation & T10 & T11 & T12 & Missed \\
SilPure, annual wafer contract & T15 & T14 & T15 & Met \\
\midrule
\textbf{Met on original commitment} & & \textbf{5 of 11} & & \textbf{5 of 11} \\
\bottomrule
\end{tabularx}
\end{table}

The two rows set in bold are the only ones on which the arms differ. The closed loop saved
the \'Electis series B, which slipped in arm A, and replanned the CompoDis S1 cover, which
arm A met on its original date. Every other difference between the arms is a replanning that
re-dated a milestone arm A also delivered late: an announced slip rather than a discovered
one, which is not nothing and is not a save.

Under this definition of truth, five of the twenty-one confident arm A forecasts of
Table~\ref{tab:app_confident} are followed by a missed original commitment within four weeks,
at T8 and T9 on NovaFab, T10 and T15 on AvioSys, and T12 on Orbitalys, against none under the
registry's current-deadline reading.

\subsection*{Arm A Rescored After Repair}
\label{app:repair_tables}

Tables~\ref{tab:app_repair_before} to \ref{tab:app_repair_after} give the complete scoring
behind Table~\ref{tab:repair_scores}. All three score the same 120 forecasts against the
original-commitment target, censored points excluded, so the observation counts are identical
across the three and only the model differs. The first is the frozen campaign as run, the
second applies the structural repair of Equation~\ref{eq:completion_repaired} alone, and the
third adds the progress-uncertainty term.

\begin{table}[htbp]
\centering
\caption{Arm A before repair, original-commitment target}
\label{tab:app_repair_before}
\small
\begin{tabularx}{\linewidth}{l *{4}{>{\centering\arraybackslash}X}}
\toprule
\textbf{Quantity} & \textbf{h = 1} & \textbf{h = 2} & \textbf{h = 3} & \textbf{h = 4} \\
\midrule
Observations & 112 & 104 & 96 & 90 \\
Positive outcomes & 7 & 14 & 20 & 26 \\
Base rate & 0.062 & 0.135 & 0.208 & 0.289 \\
\midrule
Brier score & 0.1335 & 0.1969 & 0.2728 & 0.3190 \\
Brier skill & $-1.278$ & $-0.690$ & $-0.654$ & $-0.553$ \\
AUC & 0.736 & 0.636 & 0.641 & 0.654 \\
AUC 95\% CI & [0.49, 0.93] & [0.37, 0.87] & [0.39, 0.85] & [0.42, 0.85] \\
PR-AUC & 0.163 & 0.205 & 0.266 & 0.353 \\
\bottomrule
\end{tabularx}
\end{table}

\begin{table}[htbp]
\centering
\caption{Arm A after the structural repair alone}
\label{tab:app_repair_struct}
\small
\begin{tabularx}{\linewidth}{l *{4}{>{\centering\arraybackslash}X}}
\toprule
\textbf{Quantity} & \textbf{h = 1} & \textbf{h = 2} & \textbf{h = 3} & \textbf{h = 4} \\
\midrule
Brier score & 0.0965 & 0.1633 & 0.2243 & 0.2613 \\
Brier skill & $-0.646$ & $-0.401$ & $-0.360$ & $-0.272$ \\
AUC & 0.793 & 0.671 & 0.657 & 0.721 \\
AUC 95\% CI & [0.60, 0.95] & [0.51, 0.88] & [0.47, 0.87] & [0.54, 0.92] \\
PR-AUC & 0.355 & 0.305 & 0.355 & 0.484 \\
\bottomrule
\end{tabularx}
\end{table}

\begin{table}[htbp]
\centering
\caption{Arm A after the structural repair and the progress-uncertainty term}
\label{tab:app_repair_after}
\small
\begin{tabularx}{\linewidth}{l *{4}{>{\centering\arraybackslash}X}}
\toprule
\textbf{Quantity} & \textbf{h = 1} & \textbf{h = 2} & \textbf{h = 3} & \textbf{h = 4} \\
\midrule
Brier score & 0.0988 & 0.1593 & 0.2120 & 0.2499 \\
Brier skill & $-0.687$ & $-0.367$ & $-0.286$ & $-0.216$ \\
AUC & 0.729 & 0.634 & 0.602 & 0.584 \\
AUC 95\% CI & [0.56, 0.95] & [0.42, 0.83] & [0.29, 0.85] & [0.21, 0.87] \\
PR-AUC & 0.195 & 0.228 & 0.299 & 0.385 \\
\bottomrule
\end{tabularx}
\end{table}

The confidence intervals are cluster bootstraps resampling whole nodes rather than individual
node-weeks, because forecasts on one node are strongly correlated and an i.i.d.\ bootstrap
would understate the width. They are wide enough at every horizon to contain the AUC of every
other table in this group, which is why the AUC movements between the three are reported as
uninterpretable rather than as gains or losses, while the Brier skill movements, which do not
depend on the size of the positive class in the same way, are reported as measured.

\subsection*{Distribution of the Milestone Probability}
\label{app:binarity}

Table~\ref{tab:app_binarity} counts how many of the 120 milestone-miss probabilities fall at
the extremes rather than in a usable band, which is the measurement behind the
progress-uncertainty term of Section~\ref{sec:res_repair}.

\begin{table}[htbp]
\centering
\caption{Degeneracy of the milestone-miss probability, 120 forecasts}
\label{tab:app_binarity}
\small
\begin{tabularx}{\linewidth}{>{\raggedright\arraybackslash}p{5.2cm} >{\centering\arraybackslash}X >{\centering\arraybackslash}X >{\centering\arraybackslash}X}
\toprule
\textbf{Model} & \textbf{Exactly 0 or 1} & \textbf{In $(0.05, 0.95)$} &
\textbf{Forecasts $\ge 0.5$ at 4 weeks} \\
\midrule
Frozen campaign as run & 102 (85.0\%) & 14 & 21 \\
Structural repair alone & 89 (74.2\%) & 10 & 10 \\
Structural repair with progress uncertainty & 86 (71.7\%) & 17 & 10 \\
\bottomrule
\end{tabularx}
\end{table}

The structural repair alone reduces the count of extreme values but not the usable band: it
removes eleven false certainties without producing gradations. The progress-uncertainty term
is what widens the usable band, from 10 forecasts to 17. Neither is sufficient, and $71.7\%$
of the probabilities remain degenerate.
